\section{One uncalibrated law at an oblique contact}
\label{sec:v72-one-law}

The two-offset observation separates unequal actions without using their
first derivatives.  At an oblique contact those derivatives are already
fixed by the geometric marks.  They remove both the multiplicative
ambiguity of a rank-one factorization and the unknown time scale.  This
gives a second observation theorem for the same smooth inverse, without
changing the two-offset theorem at normal incidence.

\subsection{An anchored factorization and its scale}

Let $I=[-R,R]$.  Consider a positive $C^M$ function on a neighborhood of
$I^2$, $M\ge2$, of the form
\begin{equation}\label{eq:v72-model}
 f(u,v)=Z^{-1}B_-(u)B_+(v)\{d-A(u)-C(v)\},
\end{equation}
where $d>0$, $Z>0$, $B_\pm>0$, $B_\pm(0)=1$, and
\begin{equation}\label{eq:v72-slopes}
 A(0)=C(0)=0,\qquad A'(0)=-p,\qquad C'(0)=p,\qquad p\ne0.
\end{equation}
The signed number $p$ is known; $d$, both amplitudes and $Z$ are unknown.
Assume $d-A-C>0$ on $I^2$.  No equality of $A$ and $C$ is imposed.
Multiplication of $f$ by an unknown positive constant is permitted.
Define
\begin{equation}\label{eq:v72-cross-ratio}
 \mathcal R_f(u,v)=\frac{f(u,v)f(0,0)}{f(u,0)f(0,v)},\qquad
 \kappa_f=\partial_u\partial_v\mathcal R_f(0,0).
\end{equation}
For a continuous positive density, its law determines its values on the
square, and thus these expressions whenever the indicated derivatives
exist.  Finite samples are treated below, not identified with pointwise
density observations.

\begin{proposition}[Identification of the unknown offset]
\label{prop:v72-clock}
For \eqref{eq:v72-model}--\eqref{eq:v72-slopes},
\begin{equation}\label{eq:v72-clock}
 \mathcal R_f=1-UV,\qquad U=\frac{A}{d-A},\quad V=\frac{C}{d-C},
 \qquad \kappa_f=\frac{p^2}{d^2}>0,\qquad
 d=\frac{|p|}{\sqrt{\kappa_f}}.
\end{equation}
Consequently the offset is determined by the single conditional law and
the marked momentum, without a count or an amplitude model.
\end{proposition}
\begin{proof}
The axis residuals $d-A(u)$ and $d-C(v)$ are positive.  Cancellation of
$Z$ and the two amplitudes gives
\[
 \mathcal R_f(u,v)=
 \frac{d\{d-A(u)-C(v)\}}{(d-A(u))(d-C(v))}
 =1-\frac{A(u)C(v)}{(d-A(u))(d-C(v))}.
\]
Since $U(0)=V(0)=0$, $U'(0)=-p/d$ and $V'(0)=p/d$,
$\kappa_f=-U'(0)V'(0)=p^2/d^2$.  The positive choice of $d$ fixes the
square-root branch.  Equivalently $\kappa_f=\partial_{uv}\log f(0,0)$,
because $\mathcal R_f$ equals one on both axes.  This identity is
insensitive to the unknown normalization or separate endpoint factors.
\end{proof}

The rank-one identity alone would determine $U,V$ only up to
$(U,V)\mapsto(\lambda U,\lambda^{-1}V)$.  More precisely, if neither
factor is identically zero and two pairs give the same product, evaluate
at a point where the second factor is nonzero to identify the first
factor up to a scalar, and then identify the second.  The specified first
derivatives in \eqref{eq:v72-slopes}, after \eqref{eq:v72-clock} has
fixed $d$, force $\lambda=1$.  The following formulas implement that
argument without differentiating a whole recovered function one extra
time.

Choose one fixed $a\in I\setminus\{0\}$ sufficiently close to zero
that $A(a)C(a)\ne0$.  Such an $a$ exists by \eqref{eq:v72-slopes}.
Using the recovered $d$, put
\begin{align}
 v_a(f)&=\frac d p\,\partial_u\mathcal R_f(0,a),&
 u_a(f)&=-\frac d p\,\partial_v\mathcal R_f(a,0),
                      \label{eq:v72-anchors}\\
 U_f(u)&=\frac{1-\mathcal R_f(u,a)}{v_a(f)},&
 V_f(v)&=\frac{1-\mathcal R_f(a,v)}{u_a(f)}.
                      \label{eq:v72-factors}
\end{align}
The quantities $u_a,v_a$ are signed scalars.  They are not positive
square roots and must not be replaced by their absolute values.

\begin{theorem}[One-law separation with an unreported offset]
\label{thm:v72-separation}
The single function $f$ in \eqref{eq:v72-model} and the known $p\ne0$
determine $d,A,C,B_-,B_+$ uniquely on $I$, with
\begin{align}
 A_f(u)&=\frac{dU_f(u)}{1+U_f(u)},&
 C_f(v)&=\frac{dV_f(v)}{1+V_f(v)},\label{eq:v72-actions}\\
 B_{-,f}(u)&=\frac{f(u,0)}{f(0,0)}(1+U_f(u)),&
 B_{+,f}(v)&=\frac{f(0,v)}{f(0,0)}(1+V_f(v)).
                         \label{eq:v72-amplitudes}
\end{align}
For $M\ge2$ this inverse is locally Lipschitz from $C^M(I^2)$ to
$\mathbb R\times C^M(I)^4$ on classes with uniform bounds for the
$C^M$ norms and positive lower bounds for $f$, $\kappa_f$,
$|u_a(f)|$, $|v_a(f)|$, $1+U_f$, and $1+V_f$.  The formulas extend to
an open neighborhood of each such bounded class, even when a perturbed
function does not have representation \eqref{eq:v72-model}.
\end{theorem}
\begin{proof}
Differentiating $1-UV$ at $(0,a)$ and $(a,0)$ gives
$v_a(f)=V(a)$ and $u_a(f)=U(a)$.  Thus \eqref{eq:v72-factors} recovers
the actual factors on all of $I$, including their zeros.  Solving the
definitions of $U,V$ gives \eqref{eq:v72-actions}.  Both final
denominators are positive, since $1+U=d/(d-A)$ and
$1+V=d/(d-C)$.  The two axis identities in \eqref{eq:v72-model} give
\eqref{eq:v72-amplitudes}.  This also proves uniqueness: any other
representation satisfying the same slopes gives the same offset,
factors, actions and normalized amplitudes.

For stability, restriction to a fixed slice is bounded in $C^M$.
Products and reciprocal functions are locally Lipschitz on bounded
$C^M$ sets with positive denominator floors.  Apply these facts first to
$\mathcal R_f$.  Evaluation of its mixed derivative at the single point
$(0,0)$ is a bounded scalar functional when $M\ge2$.  The function
$x\mapsto |p|x^{-1/2}$ is Lipschitz on an interval with a positive lower
endpoint, so $d$ is locally Lipschitz.  The two first-derivative point
evaluations in \eqref{eq:v72-anchors} are bounded scalar functionals as
well.  The stipulated nonzero anchor bounds permit division in
\eqref{eq:v72-factors}; fixed slices give $C^M$ functions, with no loss
of one derivative.  Finally apply the product and reciprocal bounds to
\eqref{eq:v72-actions}--\eqref{eq:v72-amplitudes}.  For small enough
perturbations all floors persist, possibly halved.  No assertion that the
extended inverse produces an actual billiard is made.
\end{proof}

Uniform anchors follow from geometric priors rather than an assumed
conditioning bound on a measured inverse.  Suppose $|p|\ge p_*>0$,
$d_-\le d\le d_+$, and the $C^2$ norms of $A,C$ are bounded by $H_2$.
Choose a fixed $a>0$ with $H_2a\le p_*$, inside the common collar.
Taylor's formula gives $|A(a)|,|C(a)|\ge p_*a/2$.  If the positive
axis residuals are bounded above by $D_*$, then
\[
 |u_a|,|v_a|\ge \frac{p_*a}{2D_*},\qquad
 \kappa_f\ge\frac{p_*^2}{d_+^2},\qquad
 1+U,1+V\ge\frac{d_-}{D_*}.
\]
The physical relative law gives uniform amplitude and normalized-density
bounds on a common positive square.  These observations supply all the
floors used in Theorem~\ref{thm:v72-separation}.  Constants need not
remain bounded as $p_*$ or $a$ tends to zero.

\subsection{Smooth contact determination}

We return to the fixed marked polygon of
Section~\ref{sec:v65-periodic-inverse}.  Call it \emph{oblique} if
$p_b\ne0$ at every physical phase.  At each phase select one positive
offset $d_b$, fixed over the returning lengths and the independent
preparations, but not included in the reconstruction data.  The offsets
may differ between phases.  They lie in a stipulated positive physical
window; the contact gate and all event tags retain their exact meanings.
Thus an unreported offset is a fixed nuisance parameter, not a random
offset, an incorrectly recorded success tag, or an uncertain endpoint
gate.  No new clock-dependent observation is being supplied.

\begin{theorem}[Smooth rigidity from one uncalibrated law per phase]
\label{thm:v72-rigidity}
For an oblique clear nongrazing marked polygon with distinct physical
contacts and positive-curvature $C^\infty$ graphs, the collection
\[
 \bigl(f_{b,d_b}:b=0,\ldots,r-1\bigr)
\]
of single-offset limiting endpoint laws determines all the offsets
$d_b$ and the complete smooth reflecting contact germs.  Neither the
offsets, curvature, flux amplitudes, nor acceptance probabilities are
supplied.  No symmetry, analyticity or closeness between the candidates
is required.

On a bounded positive marked class as in
Theorem~\ref{thm:v68-smooth-stability}, with $m\ge3$, uniform obliquity,
offsets in $[d_-,d_+]$, and common positive-residual margins, there are
$R,C,\eta_0>0$ such that
\begin{equation}\label{eq:v72-profile-stability}
 \|F-G\|_{C^0([-R,R])}+\|d-\widetilde d\|_{\ell^\infty}
 \le C\max_b\|f^F_{b,d_b}-f^G_{b,\widetilde d_b}\|_{C^m([-R,R]^2)}
\end{equation}
when the maximum on the right is at most $\eta_0$.  These are realizable
law pairs; their normalizing constants need not coincide.
\end{theorem}
\begin{proof}
Lemma~\ref{lem:v65-signed-jacobi} supplies the actual physical law
\eqref{eq:v72-model}, with $A_b=S_b^--p_bu$ and $C_b=S_b^++p_bv$.
The half-line actions are anchored, so their physical slopes are exactly
\eqref{eq:v72-slopes}.  Theorem~\ref{thm:v72-separation} therefore
recovers $d_b$ and both actions at every phase.  Add the known $p_bu$
to $A_b$ to obtain the future action map $\mathcal S(F)$.
The action version of Theorem~\ref{thm:v68-smooth-rigidity} identifies
the actual contact germs.  In particular the last step still uses the
contracted-visit envelope on flat differences, not just equality of
Taylor series.

On the bounded class choose the common collar and fixed anchor using
the preceding uniform bounds.  The inverse in
Theorem~\ref{thm:v72-separation} gives $C^m$ control of both actions
and scalar control of the offsets from the displayed law distance.
Apply the action estimate \eqref{eq:v68-real-profile-stability}; its
polynomial alignment has already been proved before the weighted norm
is used.  This gives \eqref{eq:v72-profile-stability}.  There is no extra
derivative loss from clock recovery, since $m\ge3$ and the point
functional in \eqref{eq:v72-clock} has order only two.
\end{proof}

\begin{lemma}[Normal incidence and distinct physical contacts]
\label{lem:v72-obliquity}
A clear nongrazing periodic billiard polygon with $r\ge3$ distinct
physical contacts is oblique.  If a periodic orbit with distinct
physical contacts has a normal reflection, it is the ordinary
self-retracing two-contact orbit, with zero closing translation.
\end{lemma}
\begin{proof}
At a normal reflection the outgoing ray is the incoming ray traversed
backwards.  Its first collision is therefore exactly the preceding
lifted contact: the incoming open flight was clear.  Consequently the
preceding and following physical contacts coincide.  For $r\ge3$ these
are different cyclic positions, contradicting distinctness.  The only
remaining case is $r=2$.  Here the equality of the preceding and
following lifted, not merely physical, contact forces zero closing
translation.  The two flights traverse the same segment in opposite
directions; periodicity forces normal reflection at the other endpoint
as well.  This proves the assertion without a genericity assumption.
\end{proof}

Thus the one-law theorem applies to every polygon with at least three
distinct physical contacts in the main theorem, not just a generic
subset of those polygons.  At normal incidence $\kappa_f=0$, and the
new clock formula is not defined.  This is a boundary of this argument,
not an impossibility claim about all observations at a normal contact.
The general two-calibrated-offset theorem is unchanged.  For a
self-retracing channel at one \emph{known} offset, the separate symmetric
factorization of Theorem~\ref{thm:v26-density-inverse} also remains
available.  It uses a nonzero action anchor instead of a marked nonzero
first derivative and does not identify an unknown offset by
\eqref{eq:v72-clock}.

The equal-area flat family of Proposition~\ref{prop:v68-flat-family}
can be based on the actual three-contact configuration of
Proposition~\ref{prop:v64-triangle}.  Its two distinct tables
cannot have equal single-offset limiting laws at every phase, even if
the nuisance offsets are allowed to differ: equality would first force
equal offsets and then equal smooth germs by
Theorem~\ref{thm:v72-rigidity}.  Flat contact information therefore
survives the removal of the second offset and its calibration.
Conversely, identical contact collars with identical offsets still
give the exact local-completion invariance of
Proposition~\ref{prop:v71-locality}.  The new theorem changes the local
observation, not its target to unvisited boundary arcs.

\subsection{Finite preparations with joint offset recovery}

Let $\mathscr U$ be a nonempty class of pairs $(T,d)$ consisting of an
actual marked table and one fixed offset per phase.  Impose the bounded
positive contact and forward derivative hypotheses of
Section~\ref{sec:v69-sampled-contact}, replacing the two known offsets
by $d_b\in[d_-,d_+]$ and adding the obliquity floor $|p_b|\ge p_*$.
All positive-residual, clearance and free-segment margins hold uniformly
over that interval.  The finite higher graph bounds are those needed
for order $m+s$, where $m$ is the same weighted inverse order and
$s\ge1$.  The free area has a fixed upper bound.

There are now $J=r$ groups.  A preparation in group $b$ uses the actual
fixed $d_b$; it records the same exact itinerary and $Q_+$ acceptance
tag as before, followed on success by an endpoint pair with bounded
post-acceptance error $\delta$.  Densities normalized on $Q_+$ are
restricted to $Q\Subset Q_+$ without renormalization.  No additional
clock reading, calibration trial or endpoint height is used.  Preparations
are independent resets, not successive pieces of one orbit.

\begin{theorem}[Charged joint recovery of profiles and offsets]
\label{thm:v72-sampling}
On $\mathscr U$ there is a measurable estimator
$(\widehat T,\widehat d)\in\mathscr U$ with the following property.
Let $B$ preparations be made in each of the $J=r$ groups, and set
\[
 \alpha=\frac{s}{2(m+s)+2},\qquad
 \beta=\frac{\alpha\omega}{\omega+\alpha\Gamma},\qquad
 \gamma=\frac{s}{m+s+3},\qquad L_B=\log(C_*JB/\zeta),
\]
where $\omega,\Gamma,c_0,N_0$ are uniform forward constants.
Choose
\[
 N_B=P\left\lfloor
 \frac{\alpha\log(B/L_B)}{P(\omega+\alpha\Gamma)}
 \right\rfloor,\qquad
 n_B=\lfloor Bc_0e^{-\Gamma N_B}/2\rfloor.
\]
For $0<\zeta<1/2$, assume $N_B\ge N_0$, $n_B\ge2$ and
$Bc_0e^{-\Gamma N_B}\ge8\log(J/\zeta)$.  Impose also the upper
bandwidth and small-error conditions in the proof below.  Then
\begin{equation}\label{eq:v72-rate}
 \sup_{(T,d)\in\mathscr U}\mathbb P_{T,d}\!\left(
 \|F^{\widehat T}-F^T\|_{C^0([-R,R])}
 +\|\widehat d-d\|_{\ell^\infty}
 >C\{(L_B/B)^\beta+\delta^\gamma\}\right)\le2\zeta.
\end{equation}
All $rB$ preparations, including failures, are charged.  The estimator
is an existence construction on the realizable image and the rate is a
sufficient upper bound, not an efficiency or minimax assertion.
\end{theorem}
\begin{proof}
First check the forward hypotheses with the nuisance parameter included.
The offset enters the physical numerator affinely.  Its compact range,
the positive residual integral and the differentiated relative theorem
therefore give, uniformly in $(T,d)$,
\[
 \|f_{N,b,d_b}\|_{C^{m+s}(Q_+)}+\|f_{b,d_b}\|_{C^{m+s}(Q_+)}\le H,
 \qquad \|f_{N,b,d_b}-f_{b,d_b}\|_{C^m(Q_+)}\le C_0e^{-\omega N}.
\]
The same exact reference twist and area upper bound as in
\eqref{eq:v69-acceptance} give success probability at least
$c_0e^{-\Gamma N}$.  Knowledge of the actual value of $d_b$ is not
used by any of these uniform bounds.  The limiting densities also have
a common positive floor, so Theorem~\ref{thm:v72-rigidity} applies.

For $n$ accepted endpoints in each group, use the kernel estimator of
Lemma~\ref{lem:v69-density-estimation}.  With
$L_n=\log(C_1Jn/\zeta)$ and $n^{-1}\le h\le h_0$, its error relative
to the limiting density tuple, in the product $C^m(Q)$ norm, is at most
\begin{equation}\label{eq:v72-statistical-radius}
 E_{n,N,h}=C\left\{e^{-\omega N}+h^s+
 \sqrt{\frac{L_n}{nh^{2m+2}}}+\frac{L_n}{nh^{m+2}}+
 n^{-2}+\delta h^{-m-3}\right\}
\end{equation}
with probability at least $1-\zeta$.  This lemma allows bounded
adversarial readout errors after acceptance; it does not change event
tags or replace the fixed offsets by varying offsets.

For clarity, include the nuisance parameter in the realizability step.
The subset
$\{(f^T_{b,d_b}|_Q)_b:(T,d)\in\mathscr U\}$ of the finite product
$C^m(Q)^J$ is separable.  Choose a countable dense subset of that image
and one representative $(T_j,d^j)\in\mathscr U$ for each of its points.
The infimum of the distances from the fitted tuple to these points is a
measurable function.  Select the first point whose distance is within
$n^{-2}$ of that infimum.  Such a point exists, the selection is
measurable, and its distance from the true law tuple is at most
$2E_{n,N,h}+n^{-2}$.  Apply \eqref{eq:v72-profile-stability} only to
this selected realizable pair.  Require
$2E_{n,N,h}+n^{-2}\le\eta_0$; this is the small-error condition.
The output includes the representative's actual offset tuple, not an
unconstrained clock estimate attached to an unrelated table.

Take
\[
 h=\max\{(L_n/n)^{1/(2(m+s)+2)},\,
                   \delta^{1/(m+s+3)},\,n^{-1}\}\le h_0.
\]
This is the upper bandwidth condition.  Substitution into
\eqref{eq:v72-statistical-radius}, with constants enlarged in the
small-error regime, gives
$C\{e^{-\omega N}+(L_n/n)^\alpha+\delta^\gamma\}$.
The first $n_B$ successful marks are independent with their conditional
densities: independent reset records admit the Bernoulli/conditional-mark
representation used in Section~\ref{sec:v69-sampled-contact}.
The binomial lower-tail bound and a union bound show that all groups
have $n_B$ successes except on an event of probability at most $\zeta$.
If a group fails this requirement, choose one fixed representative of
$\mathscr U$; otherwise apply the preceding selection to the first
$n_B$ successes.  This defines an estimator for every record.

Finally the specified returning length satisfies
\[
 e^{-\omega N_B}\le C(L_B/B)^\beta,\qquad
 n_B\ge cB^{\omega/(\omega+\alpha\Gamma)}
             L_B^{\alpha\Gamma/(\omega+\alpha\Gamma)}.
\]
The floor changes constants only, under the stipulated lower-count
condition.  Since $n_B\le B$ and $L_{n_B}\le CL_B$, one has
$(L_{n_B}/n_B)^\alpha\le C(L_B/B)^\beta$.
Adding the two exceptional-event probabilities proves \eqref{eq:v72-rate}.
At fixed confidence all count conditions hold for large $B$; the
bandwidth and small-error conditions still require sufficiently small
readout error.  They are not suppressed for varying confidence.
\end{proof}

The two mathematical interfaces of the original theorem remain essential.
Without a differentiated \emph{relative} physical limit,
\eqref{eq:v72-clock} could not be estimated from rare finite bridges.
Without the actual envelope inverse, recovery of the action functions
would not identify flat changes of a smooth contact.  The new algebra
removes a second offset and its supplied numerical calibration; it is
not presented as a replacement for those analytic arguments.  The marks,
exact tags and gates, finite-smoothness priors and independent resetting
remain part of the observation problem.
